% ga_pipeline.tex -- standalone source for Fig. "fig:ga_pipeline"
% Build:  pdflatex ga_pipeline.tex   ->  ga_pipeline.pdf
% Use:    paste the GA-OVERLAY block printed in ga_pipeline.log into the
%         figure environment of the paper. It includes ga_pipeline.pdf at its
%         natural size and, through \figeqlink (defined in the paper's
%         preamble), adds invisible hyperlinks over the equation numbers and
%         warns if a number below no longer matches the paper.
%
% Same class and font setup as the paper, so \small, \footnotesize and the
% math fonts are identical. preview crops the page to the TikZ bounding box,
% i.e. exactly the box the inline tikzpicture occupied.
\documentclass[10pt,prd,preprintnumbers,floatfix,aps,nofootinbib,twocolumn,superscriptaddress,noshowpacs]{revtex4-1}
\usepackage[utf8]{inputenc}
\usepackage{amssymb}
\usepackage{amsmath}
\usepackage{amsfonts}
\usepackage{amsbsy}
\usepackage[dvipsnames]{xcolor}
\usepackage[T1]{fontenc}
\usepackage{newtxtext}
\usepackage{newtxmath}
\usepackage[cal=cm]{mathalfa}

\usepackage{tikz}
\usetikzlibrary{positioning,arrows.meta,fit,backgrounds,shapes.geometric,calc,tikzmark}
\definecolor{treeblue}{RGB}{48,132,196}
\definecolor{treebluedark}{RGB}{35,100,150}
\definecolor{evalorange}{RGB}{214,120,30}
% Highlight tying the offspring evaluation in step 5 to step 2.
\newcommand{\evaltag}[1]{\tikz[remember picture, baseline=(evaltag.base)]
    \node[draw=evalorange, fill=evalorange!15, rounded corners=2pt, line width=0.7pt,
          inner xsep=1.5pt, inner ysep=0.5pt, minimum height=0pt, minimum width=0pt, text width=, align=none] (evaltag) {\strut #1};}

% Equation numbers as they appear in the paper. They are printed in the figure
% itself (so the PDF is complete on its own); the paper checks them against
% its labels at every compile and warns if one has changed.
\newcommand{\eqFriedmann}{3}   % eq:friedmann
\newcommand{\eqKG}{4}          % eq:KGcosmic
\newcommand{\eqViab}{5}        % eq:viability_conditions

% Links overlaid by the paper: <tikzmark name>/<label>/<number>.
\newcommand{\eqrefmarks}{%
    bgA/eq:friedmann/\eqFriedmann,
    bgB/eq:KGcosmic/\eqKG,
    viab/eq:viability_conditions/\eqViab}

% Equation number styled like \eqref with hyperref: black parentheses, number
% in the paper's link colour (allcolors=Blue).
\newcommand{\eqlink}[2]{(\tikzmark{#1}\textcolor{Blue}{#2})}
\makeatletter\let\pct\@percentchar\makeatother  % a literal % for \typeout

\usepackage[active,tightpage]{preview}
\PreviewEnvironment{tikzpicture}
\setlength\PreviewBorder{.1pt}

\begin{document}
\begin{tikzpicture}[
    remember picture,
    node distance=5mm,
    block/.style={
        rectangle, draw=treeblue, rounded corners, line width=0.85pt,
        fill=treeblue!10!white, align=center, font=\small,
        text width=52mm,
        inner sep=3pt, minimum height=12.7mm
    },
    final/.style={
        block, fill=white, line width=1.4pt, draw=treebluedark
    },
    decision/.style={
        diamond, aspect=2, draw=treeblue, line width=0.85pt, fill=white,
        inner sep=0.5pt, font=\small
    },
    arrow/.style={-{Stealth[length=2.mm]}, thick, draw=treebluedark},
    lbl/.style={font=\footnotesize, inner sep=1.5pt}
]
    % ---------------- nodes ----------------
    \node[block] (init) {
        \textbf{1. Initialize population} \(g\!\!=\!\!1\)\\
        Random individuals \(j \!=\! 1, \dots, N_{\rm pop}\):\\
        \(\mathcal{I}_j \!\!=\!\! \{F_j(\varphi; \boldsymbol{c}_{F\!,j}), U_j(\varphi; \boldsymbol{c}_{U,j})\}\)
    };
    \node[block, below=of init, draw=evalorange, fill=evalorange!15] (eval) {
        \textbf{2. Evaluate each individual}\\
        Minimize \(\chi^2(\boldsymbol{\theta}_j, \boldsymbol{c}_j)\) subject to\\
        background~\eqlink{bgA}{\eqFriedmann}--\eqlink{bgB}{\eqKG} and viability~\eqlink{viab}{\eqViab}
    };
    \node[block, below=8mm of eval] (sel) {
        \textbf{3. Tournament selection}\\
        Compare random groups by \(\chi^2\);\\
        pick winners and losers
    };
    \node[block, below=of sel] (cross) {
        \textbf{4. Genetic operations}\\
        Crossover and mutation of the winners'\\
        \(\{F(\varphi),\, U(\varphi)\}\) produce offspring
    };
    \node[block, below=of cross] (replace) {
        \textbf{5. Next generation}\\
        Offspring, evaluated as in \evaltag{step 2},\\
        replace losers; \(g \!\!\to\!\! g\!+\!1\)
    };
    \node[decision, below=of replace] (stop) {\(g < N_{\rm gen}\)?};
    \node[final, below=8mm of stop] (termination) {
        \textbf{6. Output}\\
        Best individual \(\mathcal{I}_\star\) (lowest \(\chi^2\))\\
        and its best-fit \(\{\hat{\boldsymbol{\theta}}_\star, \hat{\boldsymbol{c}}_\star\}\)
    };

    % ---------------- generation-loop panel ----------------
    \coordinate (rail)   at ([xshift=5mm]sel.east);
    \coordinate (lrail)  at ([xshift=-5mm]sel.west);    % mirrors the rail, keeps it centred
    \coordinate (padtop) at ([yshift=5mm]sel.north);    % room for the label
    \coordinate (padbot) at ([yshift=-5mm]stop.south);  % room for "no"
    \begin{scope}[on background layer]
        \node[fit=(lrail)(sel)(stop)(rail)(padtop)(padbot),
              inner xsep=2.5mm, inner ysep=0pt, line width=0.5pt,
              fill=black!5, draw=black!30, rounded corners=4pt] (loop) {};
    \end{scope}
    \node[anchor=north west, font=\small,
    text=black!75, inner sep=2.5pt] at (loop.north west) {Generation loop};
    \node[anchor=north east, font=\small,
    text=black!75, inner sep=2.5pt] at (loop.north east) {\(g = 2, \dots, N_{\rm gen}\)};

    % ---------------- arrows ----------------
    \draw[arrow] (init) -- (eval);
    \draw[arrow] (eval) -- (sel);
    \draw[arrow] (sel) -- (cross);
    \draw[arrow] (cross) -- (replace);
    \draw[arrow] (replace) -- (stop);
    \draw[arrow] (stop) -- node[lbl, right, pos=0.25] {no} (termination);
    \draw[arrow] (stop.east) -- node[lbl, above, pos=0.5] {yes}
        (stop.east -| rail) -- (rail) -- (sel.east);

    % ---------------- equation references for the main document ----------------
    % Prints (in the .log) the GA-OVERLAY block to paste into the figure
    % environment of the paper. Needs two pdflatex runs (tikzmark); recompile
    % until the numbers stop changing.
    \typeout{^^J\pct\pct\space---- GA-OVERLAY: paste from here ----^^J%
        \string\begin{overpic}[abs,unit=1pt]{ga_pipeline.pdf}}
    \foreach \mark/\lab/\num in \eqrefmarks {
        \path[overlay] let \p1 = ($(pic cs:\mark) - (current bounding box.south west)$) in
            \pgfextra{%
                \pgfmathsetmacro\refx{\x1}%
                \pgfmathsetmacro\refy{\y1}%
                \pgfmathprintnumberto[fixed, precision=2, assume math mode=true]{\refx}\refx
                \pgfmathprintnumberto[fixed, precision=2, assume math mode=true]{\refy}\refy
                \typeout{\space\space\string\figeqlink{\refx}{\refy}{\lab}{\num}\pct}%
            };
    }
    % Step-5 tag -> step-2 box: a link rectangle over the tag and a target at the
    % top-left corner of step 2 (\figsteplink and \figsteptarget, paper preamble).
    \path[overlay] let \p1 = ($(evaltag.south west) - (current bounding box.south west)$),
                       \p2 = ($(evaltag.north east) - (evaltag.south west)$),
                       \p3 = ($(eval.north west) - (current bounding box.south west)$) in
        \pgfextra{%
            \pgfmathsetmacro\tagx{\x1}\pgfmathprintnumberto[fixed, precision=2, assume math mode=true]{\tagx}\tagx
            \pgfmathsetmacro\tagy{\y1}\pgfmathprintnumberto[fixed, precision=2, assume math mode=true]{\tagy}\tagy
            \pgfmathsetmacro\tagw{\x2}\pgfmathprintnumberto[fixed, precision=2, assume math mode=true]{\tagw}\tagw
            \pgfmathsetmacro\tagh{\y2}\pgfmathprintnumberto[fixed, precision=2, assume math mode=true]{\tagh}\tagh
            \pgfmathsetmacro\tgtx{\x3}\pgfmathprintnumberto[fixed, precision=2, assume math mode=true]{\tgtx}\tgtx
            \pgfmathsetmacro\tgty{\y3}\pgfmathprintnumberto[fixed, precision=2, assume math mode=true]{\tgty}\tgty
            \typeout{\space\space\string\figsteptarget{\tgtx}{\tgty}{ga:step2}\pct}%
            \typeout{\space\space\string\figsteplink{\tagx}{\tagy}{\tagw}{\tagh}{ga:step2}\pct}%
        };
    \typeout{\string\end{overpic}^^J\pct\pct\space---- GA-OVERLAY: to here ----^^J}
\end{tikzpicture}
\end{document}
